\documentclass[11pt]{article}
\usepackage[spanish]{babel}
\usepackage[utf8]{inputenc}
\usepackage{booktabs}
\usepackage{geometry}
\geometry{margin=2.5cm}

\title{VIOG: Ponderación por Inverso de la Varianza del Output Gap}
\author{}
\date{}

\begin{document}
\maketitle

El VIOG (\emph{Variance-Weighted Inverse Output Gap}) estima la brecha del producto de Colombia
combinando cinco filtros de extracción de tendencia --Baxter-King, Christiano-Fitzgerald,
Butterworth, Hodrick-Prescott potenciado (BHP) y un filtro de Kalman/UCM equivalente al comando
\texttt{ucm} de Stata-- aplicados sobre el logaritmo del PIB trimestral desestacionalizado
(1994--2023). Cada filtro produce una brecha $g_{i,t} = \ln Y_t - \ln Y^{*}_{i,t}$, y el
ponderador de cada método se define como el inverso de la suma acumulada de sus brechas
absolutas hasta el periodo $t$, normalizada por el número de observaciones. El PIB potencial
final es el promedio ponderado $\ln Y^{*}_t = \sum_i w_{i,t}\,\ln Y^{*}_{i,t}$, de modo que los
filtros más estables (menor volatilidad histórica) reciben mayor peso relativo en cada instante.

El filtro de Kalman se especifica como un modelo de componentes no observados con tendencia
local lineal y ciclo estocástico amortiguado (periodo acotado entre 6 y 32 trimestres), estimado
sobre $\ln(\text{PIB})$ mediante máxima verosimilitud con valores iniciales anclados en el
óptimo económico (evitando el máximo degenerado de varianza cíclica nula). Las primeras cuatro
observaciones (1994) se descartan por \emph{burn-in}, replicando el tratamiento de Stata. La
Tabla~\ref{tab:filtros} resume las características y el desempeño de cada filtro sobre la serie
colombiana.

\begin{table}[h]
\centering
\caption{Comparación de filtros empleados en el VIOG (Colombia, 1994--2023)}
\label{tab:filtros}
\begin{tabular}{lccc}
\toprule
\textbf{Filtro} & \textbf{Tipo} & \textbf{Parámetro clave} & \textbf{Desv. estándar brecha} \\
\midrule
Baxter-King (BK)        & Pasa-banda           & Ciclo 6--32 trim.\ ($K=12$) & 0.028 \\
Christiano-Fitzgerald    & Pasa-banda asimétrico & Ciclo 6--32 trim.           & 0.024 \\
Butterworth (BW)         & Pasa-bajo             & Orden 8, $\omega_c=1/16$    & 0.031 \\
HP potenciado (BHP)      & Suavizado iterado     & $\lambda=1600$, 3 iter.\    & 0.026 \\
Kalman / UCM             & Estado-espacio        & Ciclo 6--32 trim.\ amort.\  & 0.022 \\
\bottomrule
\end{tabular}
\end{table}

\end{document}
